\chapter{Model Predictive Control}
\label{ch:ch21}
% Function names for pseudocode are prefixed with "Mpc" to stay unique
% across the book.
\SetKwFunction{MpcLoop}{MpcLoop}
\SetKwFunction{MpcStep}{MpcStep}
\SetKwFunction{MpcSolveQp}{SolveQP}
\SetKwFunction{MpcAdmm}{AdmmQP}

\begin{objectives}
\begin{itemize}
  \item Explain the receding-horizon principle and write down the finite-horizon optimal control problem that a model predictive controller solves at every step: double-integrator model, quadratic cost, input and velocity bounds.
  \item Turn the constraints that matter for a swarm (intruder avoidance, formation keeping, communication range) into linear inequalities, and decide when a constraint should be hard and when it should be soft with a slack variable.
  \item Condense the horizon into one quadratic program: derive the prediction matrices, write $\mat{H}$, $\vect{f}$, $\mat{G}$ and the bounds, and solve the program with a small ADMM solver.
  \item Trace the MPC loop on a worked example in which an intruder crosses the reference path, identify the active constraint and explain the resulting deviation.
  \item State what terminal costs, terminal sets and slack variables buy in terms of stability and feasibility, and recognise the failure modes: infeasible hard constraints, stale linearisations, weights in inconsistent units and horizons too short to stop.
  \item Implement an MPC in \python that follows a path while keeping a formation offset and a communication distance, with prediction-aware, uncertainty-inflated avoidance constraints.
\end{itemize}
\end{objectives}

% ---------------------------------------------------------------------
\section{Why this matters for a drone swarm}
\label{sec:ch21-motivation}

Picture four quadrotors that fly in a diamond formation along the corridor of
a warehouse. The global planner of \cref{ch:ch09} has given each of them a
time-indexed path, and the paths are free of conflicts among the four. That is
not the end of the story. Each drone still has to \emph{execute} its path: turn
a list of waypoints into accelerations that respect its limits $a_{\max}$ and
$v_{\max}$, stay within a tolerance $e_{\max}$ of its slot in the diamond,
stay within the radio range $R_{\mathrm{comm}}$ of the drone it relays for,
and, when the prediction layer of \cref{ch:ch20} reports a forklift-mounted
drone that crosses the corridor, keep at least $r_{\mathrm{safe}}$ away from
its predicted positions. All of these requirements are \emph{constraints}, and
the milestone of Week~11 of the study plan (\cref{ch:appA}) states the
problem exactly: constraints must be enforced explicitly, not repaired after
they have been violated.

The reactive methods of Part~IV do part of this job. \orca (\cref{ch:ch13})
and the dynamic window approach (\cref{ch:ch14}) choose one velocity for the
next step, and they are fast and robust. But they look one step or one fixed
time $\ttc$ ahead, they treat one constraint at a time, and they know nothing
about a formation. When the corridor narrows, a one-step method notices too
late that it should have slowed down two seconds ago. \textbf{Model
predictive control}\index{model predictive control} (MPC) closes this gap.
At every control step it optimises a whole trajectory over the next $N$ steps,
subject to the drone's dynamics and to \emph{every} constraint at once, applies
only the first input, and repeats from the newly measured state. The optimiser
sees the narrowing corridor $N$ steps ahead and brakes in time; the formation
tolerance and the radio range are inequalities in the same problem; and if an
intruder's predicted disc intersects the plan, the plan bends around it. MPC is
the standard tool of process control since the 1980s
\cite{garcia1989model} and of quadrotor control since the 2010s; its theory is
mature \cite{mayne2000constrained,rawlings2017mpc}, and its computational core,
a quadratic program, is a solved problem \cite{boyd2004convex}.

In the hybrid architecture of \cref{ch:ch24}, MPC is the tracker that sits
between the planners and the flight controller: it follows the \cbs path,
takes the \orca velocity as its reference when the safety layer is active,
and is the one component that can promise, before the fact, that the
formation and communication constraints will hold. This chapter builds it
from the ground up. \Cref{sec:ch21-intuition} explains the receding-horizon
principle, \cref{sec:ch21-problem} states the optimal control problem and the
swarm constraints, \cref{sec:ch21-algorithm} condenses it into a quadratic
program and solves it, \cref{sec:ch21-example} works through an intruder
crossing with real numbers, \cref{sec:ch21-properties} discusses stability
and feasibility honestly, and \cref{sec:ch21-variants} covers tuning,
prediction-aware constraints and a generated experiment. Everything is
backed by \code{code/ch21\_mpc.py}.

% ---------------------------------------------------------------------
\section{The idea in plain words}
\label{sec:ch21-intuition}

Suppose you drive on a winding road at night. Your headlights show the next
hundred metres. You do not plan the whole trip; you plan the next hundred
metres, taking into account the speed you have, the curve you see, the car
ahead and the limits of your brakes, and then you execute the first fraction
of a second of that plan. A moment later the headlights show a slightly
different picture, and you plan again. Everything that made the first plan
sensible is still there, and everything you could not know (the car ahead
braking, the pothole) is corrected as soon as it appears.

\begin{figure}[tb]
  \centering
  \inputfigure{ch21/idea}
  \caption{The receding-horizon principle. At time $t$ the controller predicts
  the next $N$ positions from the measured state $\state_t$, optimises the $N$
  inputs so that the plan tracks the reference and stays outside the predicted
  intruder disc, and applies only $\ctrl_0$. One step later the state is
  measured again, the horizon has moved by one step, and a new plan is
  optimised; the old plan, shifted by one step, serves as the starting guess
  (warm start).}
  \label{fig:ch21-idea}
\end{figure}

MPC is this behaviour written as an algorithm (\cref{fig:ch21-idea}). The
headlights are the \textbf{prediction horizon}\index{prediction horizon} of
$N$ steps of length $\dt$; the road you see is the reference trajectory and
the predicted obstacles for the times $t+\dt,\dots,t+N\dt$; the plan is a
sequence of $N$ inputs $\ctrl_0,\dots,\ctrl_{N-1}$ chosen to minimise a cost
(distance to the reference, effort) subject to the dynamics and the
constraints; ``executing the first fraction'' is applying $\ctrl_0$ and
discarding the rest; and ``planning again'' is the \textbf{receding
horizon}\index{receding horizon}: the window moves forward by one step at
every step. Three points deserve emphasis.

\begin{itemize}
  \item \emph{The model is inside the optimiser.} The plan is not a curve
  drawn on the map; it is a sequence of states that the double integrator can
  actually reach with the chosen inputs. A plan can therefore not ask for more
  acceleration than the drone has.
  \item \emph{Constraints are part of the problem, not a filter after it.}
  If the plan satisfies $\norm{\pos_k-\hat{\vect{o}}_k}\ge r_{\mathrm{safe}}$
  for every predicted intruder position $\hat{\vect{o}}_k$, and the model is
  right, the next real state will satisfy it too. This is what
  ``explicitly enforced'' means.
  \item \emph{Feedback comes from re-measuring.} Because only the first input
  is applied and the next problem starts from the measured state, model errors,
  wind and a mispredicted intruder are corrected at the next step. MPC is a
  feedback controller that happens to compute its feedback by optimisation.
\end{itemize}

The reader will recognise the pattern. Windowed prioritized planning
(\cref{ch:ch08}) plans a window of the path and re-plans as the window moves;
\dstarlite (\cref{ch:ch05}) repairs a plan when the world changes. MPC does
the same thing in continuous state space with dynamics, and it does it at the
control rate, ten to a hundred times per second.

\begin{keyidea}
At every step, solve a small optimisation problem over the next $N$ steps that
contains the dynamics and every constraint you care about; apply the first
input, throw the rest away, and repeat from the measured state. With a linear
model, a quadratic cost and linear (or linearised) constraints, that problem
is a quadratic program that a few dozen matrix--vector products solve.
\end{keyidea}

% ---------------------------------------------------------------------
\section{Problem statement and notation}
\label{sec:ch21-problem}

\subsection{The drone model}

We use the discrete-time double integrator of \cref{ch:ch02}. The state is
$\state=(\pos,\vel)\in\R^{2n}$ with $n\in\set{2,3}$, the input is the
acceleration $\ctrl=\acc\in\R^n$, and with a zero-order hold over a step of
length $\dt$ the exact update is $\state_{k+1}=\mat{A}\state_k+\mat{B}\ctrl_k$.
In the plane,
\begin{equation}
\mat{A}=\begin{pmatrix}1&0&\dt&0\\0&1&0&\dt\\0&0&1&0\\0&0&0&1\end{pmatrix},\qquad
\mat{B}=\begin{pmatrix}\tfrac12\dt^2&0\\0&\tfrac12\dt^2\\\dt&0\\0&\dt\end{pmatrix},
\label{eq:ch21-model}
\end{equation}
and in three dimensions the same blocks appear with $3\times3$ identities. The
physical limits $\norm{\acc}\le a_{\max}$ and $\norm{\vel}\le v_{\max}$ are
discs, which are convex but not linear. Throughout this chapter we replace
them by the per-axis boxes $\abs{u_{k,i}}\le a_{\max}$ and $\abs{v_{k,i}}\le
v_{\max}$, which are linear; if you need the disc exactly, use a box of
half-width $a_{\max}/\sqrt{2}$ (which lies inside the disc) or the inscribed
polygon of \cref{prop:ch21-polygon} below. Time is indexed by $k$
\emph{within} the horizon ($k=0$ is now, $k=N$ is the end of the window) and
by $t$ in the closed loop.

\subsection{The finite-horizon optimal control problem}

\begin{definition}[Finite-horizon optimal control problem]
\label{def:ch21-ocp}
Given the current state $\state_0=\state_t$, reference states
$\state_k^{\mathrm{ref}}=(\pos^{\mathrm{ref}}(t+k\dt),\vel^{\mathrm{ref}}(t+k\dt))$
for $k=1,\dots,N$, symmetric weight matrices $\mat{Q}\succeq\mat{0}$,
$\mat{R}\succ\mat{0}$ and $\mat{P}\succeq\mat{0}$, the problem is
\begin{equation}
\begin{aligned}
\min_{\ctrl_0,\dots,\ctrl_{N-1}}\quad & J=\sum_{k=1}^{N-1}(\state_k-\state_k^{\mathrm{ref}})\T\mat{Q}\,(\state_k-\state_k^{\mathrm{ref}})
+\sum_{k=0}^{N-1}\ctrl_k\T\mat{R}\,\ctrl_k
+(\state_N-\state_N^{\mathrm{ref}})\T\mat{P}\,(\state_N-\state_N^{\mathrm{ref}})\\
\st & \state_{k+1}=\mat{A}\state_k+\mat{B}\ctrl_k,\qquad k=0,\dots,N-1,\\
& \abs{u_{k,i}}\le a_{\max},\quad \abs{v_{k,i}}\le v_{\max},\quad \pos_k\in\mathcal{X},\qquad k=1,\dots,N,\\
& \text{the swarm constraints of \cref{sec:ch21-constraints}.}
\end{aligned}
\label{eq:ch21-ocp}
\end{equation}
The \textbf{stage cost}\index{stage cost} weighs the tracking error with
$\mat{Q}$ and the effort with $\mat{R}$; the \textbf{terminal
cost}\index{terminal cost} $\mat{P}$ weighs the last predicted state;
$\mathcal{X}$ is an optional polytope (a geofence). The \textbf{MPC control
law}\index{model predictive control!control law} applies the first element
of the minimiser: $\ctrl_t=\ctrl_0^{*}(\state_t)$.
\end{definition}

A warning about letters. The matrices $\mat{Q}$ and $\mat{R}$ here are
\emph{cost weights} chosen by the designer: $\mat{Q}$ says how much a metre
of position error or a metre per second of velocity error costs, and
$\mat{R}$ how much a metre per second squared of acceleration costs. They
have nothing to do with the process and measurement noise covariances that the
Kalman filter of \cref{ch:ch18} also calls $\mat{Q}$ and $\mat{R}$; the
letters coincide by tradition, not by meaning. A natural choice is
$\mat{Q}=\diag(q_p\mat{I},q_v\mat{I})$ with $q_p\gg q_v$ and
$\mat{R}=r\mat{I}$, and the terminal weight $\mat{P}$ is best taken from the
Riccati equation (\cref{sec:ch21-properties}). Increasing $\mat{Q}$ relative
to $\mat{R}$ makes the drone track harder and use more acceleration; the
absolute scale of the cost is irrelevant.

Why a quadratic cost? Three reasons. It penalises large errors more than small
ones, which matches what we want; with the linear model it makes the problem
convex, so the solver finds the global minimum every time; and without
constraints its solution is the classical linear-quadratic regulator (LQR),
whose properties are completely understood \cite{rawlings2017mpc}. MPC is LQR
plus constraints plus a moving reference, and the constraints are the point.

\subsection{Constraints that matter for the swarm}
\label{sec:ch21-constraints}

\Cref{tab:ch21-constraints} lists the constraints a swarm drone has to
respect. The input and velocity boxes are linear and always satisfiable. The
other three involve other agents, and each needs a small piece of geometry
before it fits into a convex program.

\begin{table}[tb]
\centering\small
\caption{The constraints of a swarm drone, whether they are convex in the
drone's own positions, how they enter the quadratic program, and whether they
are usually hard or soft. The predicted positions of others
($\hat{\vect{o}}_k$, $\pos_k^{j}$) are data from the prediction layer or from
the neighbour's communicated plan.}
\label{tab:ch21-constraints}
\begin{tabular}{@{}lllll@{}}
\toprule
Constraint & Form & Convex & In the QP & Hard/soft\\
\midrule
Input limit & $\abs{u_{k,i}}\le a_{\max}$ & yes & box on $\mat{U}$ & hard\\
Velocity limit & $\abs{v_{k,i}}\le v_{\max}$ & yes & box on $\mat{S}_u^{v}\mat{U}$ & hard\\
Avoidance & $\norm{\pos_k-\hat{\vect{o}}_k}\ge r_k$ & no & one half-plane per step (\cref{def:ch21-halfplane}) & soft\\
Formation & $\norm{\pos_k^{i}-\pos_k^{j}-\vect{d}_{ij}}\le e_{\max}$ & yes & inscribed polygon (\cref{prop:ch21-polygon}) & soft\\
Communication & $\norm{\pos_k^{i}-\pos_k^{j}}\le R_{\mathrm{comm}}$ & yes & inscribed polygon & soft\\
\bottomrule
\end{tabular}
\end{table}

\paragraph{(1) Obstacle and intruder avoidance.}
The prediction layer provides, for every step $k$ of the horizon, a predicted
intruder position $\hat{\vect{o}}_k$ (\cref{ch:ch18,ch:ch20}); static
obstacles are the special case $\hat{\vect{o}}_k=\vect{o}$. With the radii of
the drone and the intruder and a margin folded into $r_k$ (in the simplest
case $r_k=r_{\mathrm{safe}}$ for all $k$), the requirement is
$\norm{\pos_k-\hat{\vect{o}}_k}\ge r_k$. This set, the outside of a disc, is
\emph{not convex}: the midpoint of two safe positions on opposite sides of the
intruder is the intruder. A quadratic program cannot express it. The standard
remedy replaces the disc by a half-plane that touches it
(\cref{fig:ch21-halfplane}).

\begin{definition}[Linearised avoidance constraint]
\label{def:ch21-halfplane}
Let $\bar{\pos}_k$ be a guess of the drone's position at step $k$ (the
\textbf{linearisation point}\index{linearisation point}) with
$\bar{\pos}_k\ne\hat{\vect{o}}_k$, and let
$\vect{n}_k=(\bar{\pos}_k-\hat{\vect{o}}_k)/\norm{\bar{\pos}_k-\hat{\vect{o}}_k}$
be the unit vector from the predicted intruder towards the guess. The
\textbf{linearised avoidance constraint}\index{avoidance constraint!linearised}
at step $k$ is the half-plane
\begin{equation}
\vect{n}_k\T(\pos_k-\hat{\vect{o}}_k)\ge r_k .
\label{eq:ch21-halfplane}
\end{equation}
\end{definition}

\begin{figure}[tb]
  \centering
  \inputfigure{ch21/halfplane}
  \caption{Convexifying the avoidance constraint at one step $k$. The
  non-convex requirement is to stay outside the disc of radius $r_k$ around
  the predicted intruder position $\hat{\vect{o}}_k$. The unit vector
  $\vect{n}_k$ points from $\hat{\vect{o}}_k$ towards the previous plan
  $\bar{\pos}_k$; the half-plane $\vect{n}_k\T(\pos_k-\hat{\vect{o}}_k)\ge r_k$
  (hatched) touches the disc and contains none of it, so every position that
  satisfies it is safe. The price is the region to the left, which is
  collision-free but excluded. The new plan point $\pos_k$ lands on the
  boundary when the constraint is active.}
  \label{fig:ch21-halfplane}
\end{figure}

\begin{proposition}[The half-plane is conservative]
\label{prop:ch21-halfplane}
Every $\pos_k$ that satisfies \cref{eq:ch21-halfplane} satisfies
$\norm{\pos_k-\hat{\vect{o}}_k}\ge r_k$.
\end{proposition}
\begin{proof}
By the Cauchy--Schwarz inequality and $\norm{\vect{n}_k}=1$,
$\norm{\pos_k-\hat{\vect{o}}_k}=\norm{\vect{n}_k}\norm{\pos_k-\hat{\vect{o}}_k}\ge\vect{n}_k\T(\pos_k-\hat{\vect{o}}_k)\ge r_k$. \qedhere
\end{proof}

The half-plane is the tangent to the disc at the point closest to the guess,
and it excludes more than the disc: everything ``behind'' the disc as seen from
$\bar{\pos}_k$. This is harmless when the guess is on the correct side of the
intruder and harmful when it is not, which is why the choice of $\bar{\pos}_k$
matters. The MPC loop uses the previous plan shifted by one step
(\cref{sec:ch21-loop}); at the very first solve, or after a failure, it uses
the reference. An alternative with the same flavour is a \textbf{convex
corridor}\index{convex corridor}: a planner (for instance the path of
\cref{ch:ch04} inflated as in \cref{ch:ch02}) provides for every step $k$ a
convex polygon $\mathcal{C}_k$ of free space, and the constraint
$\pos_k\in\mathcal{C}_k$ is a handful of linear inequalities. Corridors are
the method of choice for static clutter; half-planes are the method of choice
for a moving intruder, whose predicted disc moves with $k$.

\paragraph{(2) Formation keeping.}
Drone $i$ should stay near its slot relative to drone $j$: with the desired
offset $\vect{d}_{ij}$ and the tolerance $e_{\max}$,
$\norm{\pos_k^{i}-\pos_k^{j}-\vect{d}_{ij}}\le e_{\max}$. This is the
\emph{inside} of a disc of radius $e_{\max}$ around the moving centre
$\vect{c}_k=\pos_k^{j}+\vect{d}_{ij}$ and it is convex; when drone $j$'s plan
(or its reference) is known and communicated, $\vect{c}_k$ is data and the
constraint is convex in $\pos_k^{i}$ alone. A disc is still not linear, and
we approximate it from inside by a polygon.

\begin{proposition}[Inscribed polygon]
\label{prop:ch21-polygon}
Let $\vect{n}_1,\dots,\vect{n}_M$ be the unit vectors at angles $2\pi m/M$ and
$e>0$. If $\vect{n}_m\T(\pos-\vect{c})\le e\cos(\pi/M)$ for all $m$, then
$\norm{\pos-\vect{c}}\le e$.
\end{proposition}
\begin{proof}
The $M$ inequalities describe the regular $M$-gon whose apothem (distance from
the centre to a side) is $e\cos(\pi/M)$; its vertices lie at distance $e$ from
$\vect{c}$, on the circle. A convex polygon is the convex hull of its
vertices, all of which lie in the disc, so the polygon lies in the disc. \qedhere
\end{proof}

With $M=8$ the polygon loses $1-\cos(\pi/8)\approx7.6\,\%$ of the radius at the
flat sides, with $M=16$ only $1.9\,\%$. In three dimensions the same idea uses
a polyhedron; the book's code handles the planar case with $M=8$.

\paragraph{(3) Communication range.}
Drone $i$ must stay within radio range of its relay $j$:
$\norm{\pos_k^{i}-\pos_k^{j}}\le R_{\mathrm{comm}}$. This is the formation
constraint with $\vect{d}_{ij}=\vect{0}$ and $e_{\max}=R_{\mathrm{comm}}$, and it
enters the program through the same polygon. The two constraints differ in
their role, not their form: formation keeping is about a slot, the
communication range is about connectivity of the whole swarm graph
(\cref{ch:ch23}), and a drone typically has one formation partner but several
communication partners.

\paragraph{(4) Soft constraints and slack variables.}
Every constraint that involves a prediction of somebody else can become
impossible to satisfy: an intruder that appears at short range, a partner that
brakes hard, a reference that runs through a newly detected obstacle. A
\textbf{hard constraint}\index{hard constraint} that cannot be satisfied
makes the whole program infeasible and the solver returns nothing, at the worst
possible moment. A \textbf{soft constraint}\index{soft constraint} replaces
$g_k(\vect{z})\le0$ by
\begin{equation}
g_k(\vect{z})\le s_k,\qquad s_k\ge0,\qquad
\text{and adds}\quad w_1 s_k+w_2 s_k^2\quad\text{to the cost},
\label{eq:ch21-slack}
\end{equation}
with a \textbf{slack variable}\index{slack variable} $s_k$ and penalty
weights $w_1,w_2\ge0$. The program is then always feasible, and when the hard
version is feasible and $w_1$ is large enough, the soft version returns the
same solution with all slacks at zero (\cref{prop:ch21-exact}). When it is
not, the solver returns the input that violates the constraints \emph{least},
in the sense of the penalty, which is exactly what a drone should do when a
collision can no longer be ruled out: minimise the damage rather than give up.
The rule of thumb that follows is in \cref{tab:ch21-constraints}: input and
velocity bounds are hard, because they are on the drone's own variables and can
always be met; avoidance, formation and communication constraints are soft,
with weights large enough that the slack is used only in emergencies. The
linear term $w_1s_k$ is what makes the penalty \emph{exact}; the quadratic
term $w_2s_k^2$ keeps the program strictly convex and spreads a violation
over several steps instead of concentrating it in one.

% ---------------------------------------------------------------------
\section{The algorithm}
\label{sec:ch21-algorithm}

\subsection{Stacking the horizon: the condensed form}
\label{sec:ch21-condensed}

The dynamics constraint in \cref{eq:ch21-ocp} is an equality that we can
eliminate. Applying $\state_{k+1}=\mat{A}\state_k+\mat{B}\ctrl_k$ repeatedly
from $\state_0$ gives $\state_k=\mat{A}^{k}\state_0+\sum_{j=0}^{k-1}\mat{A}^{k-1-j}\mat{B}\ctrl_j$,
and stacking the $N$ predicted states into one vector,
\begin{equation}
\mat{X}=\begin{pmatrix}\state_1\\\state_2\\\vdots\\\state_N\end{pmatrix}
=\underbrace{\begin{pmatrix}\mat{A}\\\mat{A}^2\\\vdots\\\mat{A}^N\end{pmatrix}}_{\mat{S}_x}\state_0
+\underbrace{\begin{pmatrix}\mat{B}&\mat{0}&\cdots&\mat{0}\\\mat{A}\mat{B}&\mat{B}&\cdots&\mat{0}\\\vdots&\vdots&\ddots&\vdots\\\mat{A}^{N-1}\mat{B}&\mat{A}^{N-2}\mat{B}&\cdots&\mat{B}\end{pmatrix}}_{\mat{S}_u}
\underbrace{\begin{pmatrix}\ctrl_0\\\ctrl_1\\\vdots\\\ctrl_{N-1}\end{pmatrix}}_{\mat{U}} .
\label{eq:ch21-prediction}
\end{equation}
The \textbf{prediction matrices}\index{prediction matrices}
$\mat{S}_x\in\R^{2nN\times2n}$ and $\mat{S}_u\in\R^{2nN\times nN}$ depend only
on the model and the horizon, so they are computed once. The block
$(k,j)$ of $\mat{S}_u$ is $\mat{A}^{k-1-j}\mat{B}$ for $j<k$ and zero
otherwise: the input at step $j$ influences the states after it, and the
influence propagates through $\mat{A}$. This is the \textbf{condensed
form}\index{condensed form}\index{model predictive control!condensed form} of
the problem: the only decision variables are the $nN$ inputs, and the states
are affine functions of them. We write $\mat{S}_x^{p}$, $\mat{S}_x^{v}$,
$\mat{S}_u^{p}$, $\mat{S}_u^{v}$ for the rows that produce the positions and
the velocities, and $\mat{S}_{u,k}^{p}$ for the $n$ rows that produce
$\pos_k$.

\subsection{The quadratic program}
\label{sec:ch21-qp}

Let $\bar{\mat{Q}}=\diag(\mat{Q},\dots,\mat{Q},\mat{P})$ ($N$ blocks, the last
one $\mat{P}$), $\bar{\mat{R}}=\diag(\mat{R},\dots,\mat{R})$ and
$\mat{X}^{\mathrm{ref}}$ the stacked reference. Substituting
\cref{eq:ch21-prediction} into the cost of \cref{eq:ch21-ocp},
\begin{equation}
J=\mat{U}\T\bigl(\mat{S}_u\T\bar{\mat{Q}}\mat{S}_u+\bar{\mat{R}}\bigr)\mat{U}
+2\,(\mat{S}_x\state_0-\mat{X}^{\mathrm{ref}})\T\bar{\mat{Q}}\mat{S}_u\,\mat{U}+\text{const},
\label{eq:ch21-condensed-cost}
\end{equation}
a quadratic function of $\mat{U}$ (\cref{exr:ch21-condensed} asks you to
verify it). Collect the inputs and the slacks in $\vect{z}=(\mat{U},\vect{s})$.
Every constraint of \cref{sec:ch21-constraints} is linear in $\vect{z}$, and
the whole problem is the \textbf{quadratic program}\index{quadratic program}
(QP) in standard form
\begin{equation}
\min_{\vect{z}}\ \tfrac12\vect{z}\T\mat{H}\vect{z}+\vect{f}\T\vect{z}
\qquad\st\ \vect{l}\le\mat{G}\vect{z}\le\vect{u},
\label{eq:ch21-qp}
\end{equation}
with the cost matrices
\begin{equation}
\mat{H}=\begin{pmatrix}2(\mat{S}_u\T\bar{\mat{Q}}\mat{S}_u+\bar{\mat{R}})&\mat{0}\\\mat{0}&2w_2\mat{I}\end{pmatrix},\qquad
\vect{f}=\begin{pmatrix}2\mat{S}_u\T\bar{\mat{Q}}(\mat{S}_x\state_0-\mat{X}^{\mathrm{ref}})\\ w_1\vect{1}\end{pmatrix},
\label{eq:ch21-qp-cost}
\end{equation}
and the constraint rows, one block per constraint type,
\begin{equation}
\mat{G}=\begin{pmatrix}\mat{I}&\mat{0}\\ \mat{S}_u^{v}&\mat{0}\\ \mat{N}^{c}\mat{S}_u^{p}&\mat{0}\\ -\mat{N}^{o}\mat{S}_u^{p}&-\mat{I}\\ \mat{0}&\mat{I}\end{pmatrix},\qquad
\vect{l}=\begin{pmatrix}-a_{\max}\vect{1}\\ -v_{\max}\vect{1}-\mat{S}_x^{v}\state_0\\ -\infty\\ -\infty\\ \vect{0}\end{pmatrix},\qquad
\vect{u}=\begin{pmatrix}a_{\max}\vect{1}\\ v_{\max}\vect{1}-\mat{S}_x^{v}\state_0\\ \vect{c}\\ \vect{b}\\ +\infty\end{pmatrix}.
\label{eq:ch21-qp-constraints}
\end{equation}
Read the blocks from top to bottom. The input box is a bound on $\mat{U}$
itself. The velocity box is a bound on the velocity rows of the prediction,
$\mat{S}_u^{v}\mat{U}$, shifted by the free response $\mat{S}_x^{v}\state_0$.
The keep-in block has one row per step $k$ and direction $m$ of the polygon:
$\mat{N}^{c}$ stacks the row vectors $\vect{n}_m\T$ applied to the position
block of step $k$, and the bound is
$c_{k,m}=e\cos(\pi/M)-\vect{n}_m\T\bigl((\mat{S}_x\state_0)^{p}_k-\vect{c}_k\bigr)$.
The avoidance block has one row per step and obstacle: the row is
$-\vect{n}_k\T\mat{S}_{u,k}^{p}$, the bound is
$b_k=\vect{n}_k\T\bigl((\mat{S}_x\state_0)^{p}_k-\hat{\vect{o}}_k\bigr)-r_k$,
and the $-\mat{I}$ next to it subtracts the slack $s_k$, so that the row
reads $\vect{n}_k\T(\pos_k-\hat{\vect{o}}_k)\ge r_k-s_k$. The last block is
$\vect{s}\ge\vect{0}$. For a hard avoidance constraint delete the slack
column and the last block. Because $\bar{\mat{R}}\succ\mat{0}$ and $w_2>0$,
$\mat{H}$ is positive definite, the program is strictly convex, and its
minimiser is unique.

The size of the program is what makes MPC practical. For the planar drone
with $N=15$ there are $30$ inputs and $15$ slacks, and with one intruder
$30+30+15+15=90$ constraint rows; \cref{sec:ch21-implementation} reports the
solve times.

\subsection{Solving the quadratic program}
\label{sec:ch21-solver}

Three families of methods solve \cref{eq:ch21-qp}
\cite{boyd2004convex,rawlings2017mpc}. \textbf{Active-set
methods}\index{quadratic program!active-set method} guess which constraints
hold with equality, solve the resulting equality-constrained problem (a
linear system), and add or drop constraints until the guess is right; they are
exact and warm-start beautifully, but each change of the active set costs a
linear solve. \textbf{Interior-point methods}\index{quadratic program!interior-point method}
follow a central path with Newton steps; they are robust and converge in a
few dozen iterations regardless of the problem, but they cannot exploit a
warm start. \textbf{First-order splitting methods}\index{quadratic program!ADMM}
such as the alternating direction method of multipliers (ADMM)
\cite{boyd2011admm} take many cheap steps, each a matrix--vector product,
and they warm-start well. The book uses a small dense ADMM solver written in
NumPy, the iteration of the OSQP solver \cite{stellato2020osqp}, for four
reasons: it is thirty lines that you can read in full; it needs no library;
the shifted previous plan is a warm start it can use; and it detects an
infeasible program instead of looping forever. For a real vehicle you would
call OSQP itself, or a general nonlinear solver such as
\code{scipy.optimize.minimize} with SLSQP when you prototype with non-linear
constraints, but the algorithm is the same.

\begin{algorithm}[htb]
\caption{ADMM for the quadratic program \cref{eq:ch21-qp} (the OSQP iteration)}
\label{alg:ch21-admm}
\KwIn{$\mat{H}$, $\vect{f}$, $\mat{G}$, $\vect{l}$, $\vect{u}$; warm start $\vect{z}$, $\vect{y}$; penalty $\rho$, regularisation $\sigma$, relaxation $\alpha$, tolerance $\eps$}
\KwOut{minimiser $\vect{z}$, multipliers $\vect{y}$, status}
\Fn{\MpcAdmm{$\mat{H},\vect{f},\mat{G},\vect{l},\vect{u},\vect{z},\vect{y}$}}{
  scale every row of $\mat{G}$ and of $\vect{l},\vect{u}$ to unit norm\label{alg:ch21-admm:scale}\;
  $\vect{w}\gets\Pi_{[\vect{l},\vect{u}]}(\mat{G}\vect{z})$; $\mat{K}\gets\mat{H}+\sigma\mat{I}+\rho\mat{G}\T\mat{G}$, factorised once\label{alg:ch21-admm:factor}\;
  \For{$i=1,2,\dots,i_{\max}$}{
    $\tilde{\vect{z}}\gets\mat{K}\inv\bigl(\sigma\vect{z}-\vect{f}+\mat{G}\T(\rho\vect{w}-\vect{y})\bigr)$\label{alg:ch21-admm:linsys}\;
    $\hat{\vect{w}}\gets\alpha\mat{G}\tilde{\vect{z}}+(1-\alpha)\vect{w}$;\quad $\vect{z}\gets\alpha\tilde{\vect{z}}+(1-\alpha)\vect{z}$\label{alg:ch21-admm:relax}\;
    $\vect{w}^{+}\gets\Pi_{[\vect{l},\vect{u}]}(\hat{\vect{w}}+\vect{y}/\rho)$\label{alg:ch21-admm:proj}\;
    $\vect{y}\gets\vect{y}+\rho(\hat{\vect{w}}-\vect{w}^{+})$;\quad $\vect{w}\gets\vect{w}^{+}$\label{alg:ch21-admm:dual}\;
    \If{$\norm{\mat{G}\vect{z}-\vect{w}}_\infty\le\eps$ \KwAnd $\norm{\mat{H}\vect{z}+\vect{f}+\mat{G}\T\vect{y}}_\infty\le\eps$}{
      polish: re-solve the KKT system on the active rows; \KwRet{$(\vect{z},\vect{y},\mathit{solved})$}\label{alg:ch21-admm:polish}\;
    }
    \If{the change of $\vect{y}$ certifies infeasibility}{\KwRet{$(\vect{z},\vect{y},\mathit{infeasible})$}\label{alg:ch21-admm:infeasible}\;}
    every $30$ iterations: rebalance $\rho$ from the ratio of the two residuals and refactorise $\mat{K}$\label{alg:ch21-admm:rho}\;
  }
  \KwRet{$(\vect{z},\vect{y},\mathit{max\_iter})$}\;
}
\end{algorithm}

\Cref{alg:ch21-admm} works with a copy $\vect{w}$ of $\mat{G}\vect{z}$ that
is forced into the box $[\vect{l},\vect{u}]$ and a vector of
\textbf{Lagrange multipliers}\index{Lagrange multiplier} $\vect{y}$ that
pushes the two copies together. Line~\ref{alg:ch21-admm:linsys} minimises
the cost plus a quadratic penalty for the disagreement between $\mat{G}\vect{z}$
and $\vect{w}$; this is a linear system with the fixed matrix $\mat{K}$, so
the factorisation (in the book's code, for a matrix with at most a few hundred
rows, simply the inverse) is computed once and every iteration costs two
matrix--vector products. Line~\ref{alg:ch21-admm:proj} clips the disagreeing
copy into the box, the cheapest operation imaginable, and
line~\ref{alg:ch21-admm:dual} updates the multipliers by the remaining
disagreement. Row scaling (line~\ref{alg:ch21-admm:scale}) matters more than
it looks: an input-bound row has entries of size one, an avoidance row has
entries of size $\dt^2$, and one penalty $\rho$ cannot serve both unless the
rows are equilibrated. The two residuals in the stopping test are the
\textbf{primal residual}\index{residual!primal} (are the constraints
satisfied?) and the \textbf{dual residual}\index{residual!dual} (is the
gradient of the cost balanced by the constraint forces?); at a solution both
are zero. ADMM converges to a modest accuracy quickly and to a high accuracy
slowly, so line~\ref{alg:ch21-admm:polish} \textbf{polishes}\index{polishing}
the answer: it reads the \textbf{active set}\index{active set} (the rows
whose multiplier is non-zero) off $\vect{y}$ and solves the
equality-constrained problem on those rows exactly, which makes the active
constraints hold to machine precision. Finally, if the program is infeasible
the multipliers grow along a direction $\delta\vect{y}$ with
$\mat{G}\T\delta\vect{y}=\vect{0}$ and
$\vect{u}\T\delta\vect{y}^{+}+\vect{l}\T\delta\vect{y}^{-}<0$, which is a
certificate that no $\vect{z}$ satisfies the bounds \cite{stellato2020osqp};
line~\ref{alg:ch21-admm:infeasible} checks for it, so that a hard constraint
that cannot be met is reported rather than iterated on forever.

The multipliers are more than a by-product. The multiplier $y_i$ of a
constraint row tells you how much the optimal cost would drop if that row were
relaxed by one unit; a row with $y_i>0$ is \textbf{active}\index{active
constraint} and shapes the solution, a row with $y_i=0$ could be deleted
without changing anything. In the worked example we read off which avoidance
step is active from exactly these numbers.

\subsection{The MPC loop}
\label{sec:ch21-loop}

\begin{algorithm}[htb]
\caption{Model predictive control loop with linearised avoidance and keep-in constraints}
\label{alg:ch21-loop}
\KwIn{model $(\mat{A},\mat{B})$, weights $(\mat{Q},\mat{R},\mat{P})$, horizon $N$, limits $a_{\max},v_{\max}$, reference $\state^{\mathrm{ref}}(\cdot)$, predicted obstacles $\hat{\vect{o}}_k$ with radii $r_k$, keep-in centres $\vect{c}_k$ with radii $e$}
\KwOut{the input $\ctrl_t$ applied at every step}
\Fn{\MpcLoop{}}{
  precompute $\mat{S}_x$, $\mat{S}_u$ and $\mat{H}$ of \cref{eq:ch21-qp-cost}; $\mat{U}^{\mathrm{prev}}\gets\vect{0}$\label{alg:ch21-loop:pre}\;
  \For{$t=0,\dt,2\dt,\dots$}{
    measure $\state_t$; fetch $\state_k^{\mathrm{ref}}$, $\hat{\vect{o}}_k$, $\vect{c}_k$ for the times $t+k\dt$, $k=1,\dots,N$\label{alg:ch21-loop:measure}\;
    $(\ctrl_t,\mat{U}^{\mathrm{prev}})\gets$ \MpcStep{$\state_t$, references, obstacles, keep-in discs, $\mat{U}^{\mathrm{prev}}$}\;
    apply $\ctrl_t$ to the drone for one step\label{alg:ch21-loop:apply}\;
  }
}
\Fn{\MpcStep{$\state_0$, references, obstacles, keep-in discs, $\mat{U}^{\mathrm{prev}}$}}{
  $\bar{\mat{U}}\gets$ $\mat{U}^{\mathrm{prev}}$ without its first input, with the last input repeated\label{alg:ch21-loop:shift}\;
  $\bar{\mat{X}}\gets\mat{S}_x\state_0+\mat{S}_u\bar{\mat{U}}$\tcp*{linearisation point}
  $\vect{f}\gets2\mat{S}_u\T\bar{\mat{Q}}(\mat{S}_x\state_0-\mat{X}^{\mathrm{ref}})$, extended by $w_1\vect{1}$ for the slacks\label{alg:ch21-loop:cost}\;
  add the input and velocity boxes to $(\mat{G},\vect{l},\vect{u})$\label{alg:ch21-loop:boxes}\;
  \ForEach{obstacle \KwAnd step $k=1,\dots,N$}{
    $\vect{n}_k\gets(\bar{\pos}_k-\hat{\vect{o}}_k)/\norm{\bar{\pos}_k-\hat{\vect{o}}_k}$ (use $\pos_k^{\mathrm{ref}}$ instead of $\bar{\pos}_k$ if the two coincide)\label{alg:ch21-loop:normal}\;
    add the row $-\vect{n}_k\T\mat{S}_{u,k}^{p}\mat{U}-s_k\le\vect{n}_k\T\bigl((\mat{S}_x\state_0)^{p}_k-\hat{\vect{o}}_k\bigr)-r_k$\label{alg:ch21-loop:avoid}\;
  }
  \ForEach{keep-in disc, step $k$ \KwAnd direction $m=1,\dots,M$}{
    add the row $\vect{n}_m\T\mat{S}_{u,k}^{p}\mat{U}\le e\cos(\pi/M)-\vect{n}_m\T\bigl((\mat{S}_x\state_0)^{p}_k-\vect{c}_k\bigr)$\label{alg:ch21-loop:keepin}\;
  }
  $(\vect{z},\vect{y},\mathit{status})\gets$ \MpcSolveQp{$\mat{H},\vect{f},\mat{G},\vect{l},\vect{u}$, warm start $(\bar{\mat{U}},\vect{0})$ and the shifted multipliers}\label{alg:ch21-loop:solve}\;
  \If{status $\ne$ solved}{
    \KwRet{$\bigl(\Pi_{[-a_{\max},a_{\max}]}(-\vel_0/\dt),\ \bar{\mat{U}}\bigr)$}\tcp*{fallback: brake}\label{alg:ch21-loop:fallback}
  }
  $\mat{U}^{*}\gets$ the first $nN$ entries of $\vect{z}$\;
  \KwRet{$(\ctrl_0^{*},\mat{U}^{*})$}\label{alg:ch21-loop:first}\;
}
\end{algorithm}

\Cref{alg:ch21-loop} is the receding-horizon principle of
\cref{fig:ch21-idea} with the quadratic program inside. Everything that
depends only on the model is computed once
(line~\ref{alg:ch21-loop:pre}): the prediction matrices and the Hessian
$\mat{H}$, which does not change from step to step because the cost weights do
not. At every step the loop measures the state and collects the data for the
$N$ future times (line~\ref{alg:ch21-loop:measure}): the reference from the
path the drone is following, the intruder predictions from the prediction
layer, the keep-in centres from the communicated plans of its partners. The
step function first builds the \textbf{warm start}\index{warm start}
(line~\ref{alg:ch21-loop:shift}): the previous plan minus the input that has
just been applied, with the last input repeated to fill the vacated slot. This
shifted plan is the linearisation point $\bar{\mat{X}}$ for the avoidance
half-planes (line~\ref{alg:ch21-loop:normal}) and the initial guess of the
solver (line~\ref{alg:ch21-loop:solve}). Only the linear term $\vect{f}$ of
the cost depends on the current state and reference
(line~\ref{alg:ch21-loop:cost}). Lines~\ref{alg:ch21-loop:boxes} to
\ref{alg:ch21-loop:keepin} assemble the constraint blocks of
\cref{eq:ch21-qp-constraints}. The solver returns either a plan, of which
line~\ref{alg:ch21-loop:first} applies the first input and keeps the rest for
the next warm start, or a failure, in which case
line~\ref{alg:ch21-loop:fallback} brakes as hard as the limits allow. The
fallback is a design decision; \cref{sec:ch21-implementation} discusses the
alternatives. With soft constraints the solver essentially never fails, which
is the strongest argument for them.

% ---------------------------------------------------------------------
\section{A worked example}
\label{sec:ch21-example}

\begin{example}[An intruder crosses the reference path]
\label{ex:ch21-crossing}
A planar drone with $a_{\max}=2$~m/s$^2$ (per axis), $v_{\max}=2$~m/s and
$\dt=0.1$~s starts at the origin with velocity $(1,0)$~m/s and follows the
reference $\pos^{\mathrm{ref}}(t)=(t,0)$~m, $\vel^{\mathrm{ref}}=(1,0)$~m/s.
An intruder starts at $(4.5,-3.0)$~m and moves with the constant velocity
$(0,0.8)$~m/s, so that it crosses the reference line at $t=3.75$~s, when the
drone would be $0.75$~m behind the crossing point. The required separation is
$r_{\mathrm{safe}}=1$~m. The prediction layer reports the intruder's
constant-velocity extrapolation exactly. The controller uses $N=15$
(a horizon of $1.5$~s), $\mat{Q}=\diag(1,1,0.1,0.1)$, $\mat{R}=0.1\,\mat{I}$,
the terminal weight from the Riccati equation,
\begin{equation*}
\mat{P}=\begin{pmatrix}9.078&0&3.166&0\\0&9.078&0&3.166\\3.166&0&2.766&0\\0&3.166&0&2.766\end{pmatrix},
\end{equation*}
and hard avoidance constraints. The simulation runs for $8$~s
($80$ steps) with the exact model \cref{eq:ch21-model} as the plant.
\end{example}

\begin{figure}[tb]
  \centering
  \inputfigure{ch21/example}
  \caption{\Cref{ex:ch21-crossing}. (a)~The reference (grey), the MPC
  trajectory (blue, dots every second), the intruder's path (red dashed,
  squares every second) and the safety disc at the moment of closest
  approach, $t=4.2$~s. The dotted curve is the plan computed at $t=3$~s: it
  slows down and dips below the reference line, on the side the intruder
  comes from, so that the intruder passes in front. (b)~Speed, separation and
  tracking error over time; the shaded band marks the steps at which an
  avoidance constraint is active. The separation touches
  $r_{\mathrm{safe}}=1$~m and never goes below it.}
  \label{fig:ch21-example}
\end{figure}

\begin{table}[tb]
\centering\small
\caption{Trace of \cref{alg:ch21-loop} on \cref{ex:ch21-crossing}, every
half second: state, applied input, separation from the intruder, tracking
error, the horizon steps $k$ whose avoidance constraint is active, and the
largest multiplier. All numbers are printed by \code{ch21\_mpc.py}.}
\label{tab:ch21-trace}
\begin{tabular}{@{}rrrrrrrrrll@{}}
\toprule
$t$ [s] & $x$ & $y$ & $v_x$ & $v_y$ & $a_x$ & $a_y$ & sep. & error & active $k$ & $\max y_i$\\
\midrule
2.0 & 2.000 & 0.000 & 1.000 & 0.000 & 0.000 & 0.000 & 2.865 & 0.000 & -- & 0\\
2.5 & 2.477 & 0.000 & 0.861 & $-0.010$ & $-0.416$ & $-0.214$ & 2.257 & 0.023 & 15 & 16.2\\
3.0 & 2.876 & $-0.045$ & 0.786 & $-0.199$ & $-0.016$ & $-0.447$ & 1.716 & 0.132 & 12, 13 & 14.2\\
3.5 & 3.264 & $-0.192$ & 0.763 & $-0.360$ & $-0.001$ & $-0.080$ & 1.236 & 0.304 & 7, 8 & 14.7\\
4.0 & 3.658 & $-0.362$ & 0.858 & $-0.254$ & 0.640 & 0.817 & 1.012 & 0.498 & 2, 3 & 14.7\\
4.5 & 4.183 & $-0.366$ & 1.223 & 0.220 & 0.318 & 0.458 & 1.016 & 0.484 & -- & 0\\
5.0 & 4.815 & $-0.222$ & 1.264 & 0.300 & $-0.150$ & $-0.140$ & 1.262 & 0.289 & -- & 0\\
6.0 & 5.983 & $-0.023$ & 1.074 & 0.092 & $-0.140$ & $-0.166$ & 2.351 & 0.029 & -- & 0\\
7.0 & 7.009 & 0.010 & 0.998 & $-0.001$ & $-0.019$ & $-0.025$ & 3.606 & 0.013 & -- & 0\\
\bottomrule
\end{tabular}
\end{table}

\Cref{fig:ch21-example} and \cref{tab:ch21-trace} show what happens. For the
first two seconds the drone is exactly on the reference and every input is
zero: the predicted intruder positions are far from every plan point, all
avoidance rows have zero multipliers, and the unconstrained minimiser of the
cost is ``do nothing''. At $t=2.1$~s the last plan point of the horizon,
$\pos_{15}$ at time $3.6$~s, would be $0.98$~m from the intruder's predicted
position at that time. The half-plane for $k=15$ becomes active with the
multiplier $16.2$: from now on the constraint, not the cost, shapes the end of
the plan. Its normal $\vect{n}_{15}$ points from the predicted intruder
(ahead and slightly above the line at $t=3.6$~s) back towards the drone and
down, so the cheapest way to satisfy it is to arrive later and lower. The
drone brakes ($a_x=-0.416$~m/s$^2$ at $t=2.5$~s) and starts to move below
the line ($a_y<0$).

As the horizon slides forward, the active step moves down the horizon: $k=12$
and $13$ at $t=3.0$~s, $k=7$ and $8$ at $t=3.5$~s, $k=2$ and $3$ at
$t=4.0$~s. The drone is now at the boundary of the safety disc and follows it:
the separation is $1.012$~m at $t=4.0$~s, reaches its minimum of exactly
$1.000$~m at $t=4.2$~s (the polished solution satisfies the active row to
machine precision, and the plant is the model), and grows again as the
intruder moves on. The intruder passes in front, $0.73$~m above the drone,
which has dipped to $y=-0.395$~m at most; the drone never stops
($v_x\ge0.763$~m/s) because it is cheaper, under $\mat{Q}$ and $\mat{R}$, to
keep moving on the far side of the disc than to wait. At $t=4.5$~s the last
avoidance row goes inactive. The drone is now $0.48$~m behind the reference
and accelerates to $1.277$~m/s at $t=4.8$~s to catch up; the velocity bound
$v_{\max}=2$~m/s is never reached because the quadratic cost prefers a
gentle recovery. By $t=6.5$~s the tracking error is below $1$~cm. Over the
whole run the RMS tracking error is $0.208$~m and the largest error
$0.532$~m, both at the price of the intruder.

Two details of the trace are worth a second look. First, two consecutive
steps are usually active at once ($12$ and $13$, $7$ and $8$): the plan
touches the disc along a short arc, not at a point, because the half-planes of
neighbouring steps are nearly parallel and the drone slides along the disc
for two steps. Second, the multipliers stay between $14$ and $16$ while the
constraint is active. This number is the price, in units of the cost, of one
more metre of clearance, and it is the number that decides how large the slack
weight $w_1$ of a soft constraint must be (\cref{prop:ch21-exact}): with
$w_1=50$ the soft controller reproduces this trajectory to the millimetre,
with $w_1=1$ it cuts into the disc (\cref{sec:ch21-tuning}).

